% ══════════════════════════════════════════════════════════════════
%  FICHE D'EXERCICES — Chapitre 1 : Généralités sur la chimie organique
%  Exercices type Baccalauréat (énoncés). Corrigés regroupés en fin d'ouvrage.
% ══════════════════════════════════════════════════════════════════
\section{Exercices type Baccalauréat}
\textit{Données utiles : masses molaires atomiques (en \si{\gram\per\mol}) :
$\ce{H}=1$ ; $\ce{C}=12$ ; $\ce{N}=14$ ; $\ce{O}=16$. Volume molaire dans les
conditions de l'expérience : $V_m = \SI{24}{\litre\per\mol}$. Les corrigés des
exercices marqués \diff{3} figurent en fin d'ouvrage.}

\rubrique{Analyse élémentaire \& détermination de formules}

\exo{}\diff{1} Un composé organique a pour formule brute $\ce{C4H10O}$. Calculer sa
masse molaire, puis les pourcentages massiques en carbone, en hydrogène et en oxygène.

\exo{}\diff{1} La combustion complète de \SI{4.4}{\gram} d'un alcane produit
\SI{13.2}{\gram} de dioxyde de carbone et \SI{7.2}{\gram} d'eau. Déterminer les
quantités de carbone et d'hydrogène dans l'échantillon, puis montrer que le composé
ne contient pas d'oxygène.

\exo{}\diff{2} Un composé organique gazeux contient en masse \SI{82.7}{\percent} de
carbone et \SI{17.3}{\percent} d'hydrogène. Sa densité de vapeur par rapport à l'air
est $d = 2{,}0$.
\begin{enumerate}[label=\arabic*)]
\item Déterminer sa masse molaire.
\item En déduire sa formule brute.
\item Écrire toutes les formules semi-développées possibles et les nommer.
\end{enumerate}

\exo{}\diff{2} On brûle complètement \SI{0.1}{\mol} d'un composé organique $A$ de
formule $\ce{C_xH_yO_z}$. On recueille \SI{17.6}{\gram} de $\ce{CO2}$ et
\SI{9.0}{\gram} d'eau. La masse molaire de $A$ est \SI{60}{\gram\per\mol}.
Déterminer la formule brute de $A$.

\exo{}\diff{2} L'analyse d'un composé organique $B$ donne, en masse :
\SI{40.0}{\percent} de C, \SI{6.7}{\percent} de H et le reste d'oxygène. Sa masse
molaire vaut \SI{60}{\gram\per\mol}.
\begin{enumerate}[label=\arabic*)]
\item Déterminer la composition centésimale complète, puis la formule brute de $B$.
\item Donner les formules semi-développées des isomères possibles comportant un
groupe acide carboxylique ou un groupe ester.
\end{enumerate}

\exo{}\diff{3} La combustion complète d'une masse $m$ d'un composé organique $C$ ne
contenant que C, H et O fournit \SI{0.44}{\gram} de $\ce{CO2}$ et \SI{0.18}{\gram}
d'eau. La masse molaire de $C$, déterminée par une autre méthode, est
\SI{60}{\gram\per\mol}, et $C$ donne un test acide positif au papier pH.
\begin{enumerate}[label=\arabic*)]
\item Sachant que $m = \SI{0.30}{\gram}$, déterminer les masses de carbone,
d'hydrogène et d'oxygène dans l'échantillon.
\item En déduire la formule brute de $C$.
\item Identifier $C$ et donner sa formule semi-développée et son nom.
\end{enumerate}

\rubrique{Isomérie}

\exo{}\diff{1} Donner toutes les formules semi-développées des isomères de formule
brute $\ce{C5H12}$ et les nommer.

\exo{}\diff{2} Pour la formule brute $\ce{C3H6O}$, écrire les formules
semi-développées correspondant : a)~à un aldéhyde ; b)~à une cétone ; c)~à un alcool
insaturé. Préciser, pour chaque couple, le type d'isomérie qui les relie (chaîne,
position ou fonction).

\exo{}\diff{2} Le but-2-ène présente une isomérie $Z/E$. Représenter les deux
stéréo-isomères, les nommer, et expliquer l'origine de cette isomérie.

\exo{}\diff{2} Identifier le(s) atome(s) de carbone asymétrique(s) dans le
butan-2-ol et dans l'acide 2-hydroxypropanoïque (acide lactique). Que peut-on en
déduire quant à l'existence d'énantiomères ?

\rubrique{Nomenclature}

\exo{}\diff{1} Nommer les composés suivants : a)~$\ce{CH3-CH(CH3)-CH2-CH3}$ ;
b)~$\ce{CH3-CH2-CH2-OH}$ ; c)~$\ce{CH3-CO-CH3}$ ; d)~$\ce{CH3-CH2-COOH}$.

\exo{}\diff{1} Écrire les formules semi-développées de : a)~2,3-diméthylbutane ;
b)~3-méthylbutan-2-ol ; c)~acide 2-méthylpropanoïque ; d)~propanoate d'éthyle.

\exo{}\diff{2} On considère le composé $\ce{CH3-CH(OH)-CH2-CHO}$.
\begin{enumerate}[label=\arabic*)]
\item Identifier les groupes fonctionnels présents.
\item Nommer le composé.
\item Ce composé est-il chiral ? Justifier.
\end{enumerate}

\rubrique{Problèmes de synthèse — type Baccalauréat}

\sujetbac[d'après Baccalauréat C, Cameroun]
On étudie un composé organique $A$ ne contenant que les éléments carbone, hydrogène
et oxygène. La combustion complète de \SI{0.58}{\gram} de $A$ produit
\SI{1.32}{\gram} de dioxyde de carbone et \SI{0.54}{\gram} d'eau. Par ailleurs, la
vaporisation de \SI{0.58}{\gram} de $A$ occupe un volume de \SI{0.24}{\litre} dans
les conditions où le volume molaire vaut \SI{24}{\litre\per\mol}.
\begin{enumerate}[label=\arabic*)]
\item Déterminer la masse molaire de $A$.
\item Calculer les masses de carbone, d'hydrogène et d'oxygène contenues dans
\SI{0.58}{\gram} de $A$, puis établir sa formule brute.
\item Écrire les formules semi-développées des isomères de $A$ ; les nommer et
préciser leur famille fonctionnelle.
\item $A$ réagit avec la 2,4-DNPH mais ne réduit pas la liqueur de Fehling.
Identifier $A$ et justifier.
\end{enumerate}

\sujetbac[Identification d'un hydrocarbure]
Un hydrocarbure gazeux $X$ a une densité de vapeur par rapport à l'air $d = 1{,}93$.
Sa combustion complète donne un volume de dioxyde de carbone égal à quatre fois le
volume de $X$ brûlé (mesurés dans les mêmes conditions).
\begin{enumerate}[label=\arabic*)]
\item Déterminer la masse molaire de $X$.
\item Montrer que la molécule de $X$ contient 4 atomes de carbone.
\item En déduire la formule brute de $X$, puis discuter les formules
semi-développées possibles (on précisera le caractère saturé ou insaturé).
\end{enumerate}
